\documentclass[10pt,a4paper]{amsart}
\usepackage[margin=19mm]{geometry}
\usepackage{amsmath,amssymb,mathtools}
\usepackage[scr=rsfs,cal=euler]{mathalfa}
\usepackage{xfrac}
\usepackage[unicode,hidelinks,bookmarks=false]{hyperref}
\newcommand{\ie}{i.e.\ }
\newcommand{\cf}{cf.\ }
\newcommand{\resp}{respectively\ }
\newcommand{\Subsection}{\subsection}
\providecommand{\nobreakdash}{\nobreak\hbox{-}}
\setlength{\emergencystretch}{2em}
\multlinegap0pt
\usepackage{amssymb}
\usepackage{mathtools}
\usepackage{csquotes}
\usepackage[shortlabels]{enumitem}
\newtheorem{lemma}{Lemma}
\newcommand{\Aut}{\mathrm{Aut}}
\newcommand{\ZZ}{\mathbb{Z}}
\title{Premaniplexes of rank $3$ and $4$ as symmetry type graphs of maniplexes}
\author{Maru\v{s}a Lek\v{s}e}
\date{Source: arXiv:2609.13090v1 (CC BY 4.0)}
\begin{document}
\maketitle

\noindent\emph{Source and licence.} Premaniplexes of rank $3$ and $4$ as symmetry type graphs of maniplexes. Version arXiv:2609.13090v1 (CC BY 4.0), distributed by arXiv under the Creative Commons Attribution 4.0 International licence, http://creativecommons.org/licenses/by/4.0/, as recorded in the arXiv licence metadata for this exact version and shown on the abstract page https://arxiv.org/abs/2609.13090v1. Full licence notice: \texttt{licenses/arxiv-2609.13090-CC-BY-4.0.txt}.\par\smallskip

\noindent\textit{Editorial scope.} Counted unit: the article's lemma rem:H1 (source0.tex lines 468--472). The article's complete original proof of it is delivered with it (source0.tex lines 474--511, one \texttt{proof} environment of 38 lines). No mathematics is written, repaired or re-derived: the statement and the proof are the article's own lines, in the article's own notation, and no retained line is substituted or re-flowed.  No further source-local context is needed: every cross-reference of every retained line resolves inside the two delivered spans.  Nothing was omitted from any retained span, so the package carries no omission record.  No retained line contains a citation, so the wrapper carries no bibliography and no bibliography entry.  No retained line contains a figure environment, an image-inclusion command or drawing code, so no image is delivered and the package declares no asset dependency.  Editorial formatting only, in addition: an AMS \texttt{amsart} preamble that restates the article's own declarations and macros used by the retained lines, namely the environments \texttt{lemma} on separate counters, together with the standard packages those lines need.  The retained environments print in this document as Lemma 1 in order of appearance, whereas the article prints the same items with the numbering of its own full document.  The complete native source of the article as distributed in the arXiv e-print for this exact version is bundled unchanged as \texttt{sources/210177/source0.tex}.
\begin{lemma}\label{rem:H1}
	Let $\Gamma$ be a finite connected graph without semiedges in which every vertex has
	valency at least $2$ and which is not a cycle. Then $\Aut(\Gamma)$ acts faithfully on
	$\mathrm{H}_1(\Gamma;\ZZ)$.
\end{lemma}

\begin{proof}
  It is enough to show that there
  exists no nonidentity $g\in\Aut(\Gamma)$ that fixes every cycle
  setwise, as well as their direction. Suppose that such a $g$ exists.
  
  %Note that because of the assumption
  %on degrees of vertices, $\Gamma$ contains a cycle. 
  Note also that $g$ cannot fix all vertices and only move some dart $x$, since then either $x$ is a loop, or the edges of $x$ and $x^g$, have the same
  endvertices. In both case $g$ does not fix all cycles together with their directions.

  We first aim to show that there exists at least one vertex that is moved by $g$ and lies on a
  cycle. Suppose that there exists no such vertex. Let $u$ be a vertex that is moved by $g$ and lies on
  no cycles. Observe that it is a cut-vertex, and that each of the at least two connected
  components of $\Gamma - u$ contains a cycle. Choose vertices $a$ and $b$ lying on cycles in two different components
  of $\Gamma-u$. Then $u$ lies on every path from $a$ to $b$. By our assumption $a$ to $b$ are fixed by $g$.
The vertex $u^g$ also lies on every path from $a$ to $b$, and $u$ and $u^g$ are at the same distance from $a$.
Therefore they
  lie at the same position on any shortest path from $a$ to $b$. This implies that $u^g=u$, which is a contradiction.
  We choose $v$ to be
  a vertex that is moved by $g$ and lies on a cycle.

  Now let $C$ be the cycle containing $v$. Since $g$ fixes the cycles
  setwise, and their directions, $C$ also contains $v^g$, and $g$ induces
  a nontrivial rotation on $C$. Let $\Gamma \setminus C$ denote the graph obtained from $\Gamma$ by
  deleting the edges of $C$. Each of its components contains a vertex of $C$.
  Now we separate two options. Either there exist a walk $W$
  in $\Gamma \setminus C$ connecting two distinct vertices in $C$, or each component of
  $\Gamma \setminus C$ contains exactly one vertex of $C$. In the first case, there exists a cycle $C'$ containing $W$
  and an arc of $C$. Since $g$ also fixes $C'$ setwise, and preserves its direction,
  it fixes that arc, and hence its endvertices, contradicting that it acts on $C$ as a nontrivial rotation.
  In the second case, each component of
  $\Gamma \setminus C$
  contains a cycle whenever it contains at least two vertices. Such a component exists because
  $\Gamma$ is not a cycle. The
  components are permuted by $g$, and none of them is fixed, hence $g$ moves a cycle.

  In both cases we arrive to a contradiction, which completes the proof.
\end{proof}

\end{document}
